\documentclass[11pt,a4paper]{article}

\usepackage[utf8]{inputenc}
\usepackage[T1]{fontenc}
\usepackage[spanish]{babel}
\usepackage{amsmath,amssymb}
\usepackage{booktabs}
\usepackage{geometry}
\usepackage{graphicx}
\usepackage{xcolor}

% --- Sustitutos ligeros de siunitx (paquete no disponible) ---
\newcommand{\km}{km}
\newcommand{\s}{s}
\newcommand{\Hz}{Hz}
\newcommand{\m}{m}
\newcommand{\num}[1]{#1}
\newcommand{\SI}[2]{#1\,#2}
\newcommand{\SIrange}[3]{#1--#2\,#3}
\usepackage[hidelinks]{hyperref}
\geometry{margin=2.5cm}

\title{Inversión cinemática del sismo de Calama 2020 ($M_w$ 6.8):\\
de la fuente puntual a la función de Green empírica}
\author{Alex Villarroel}
\date{12 de junio de 2026}

\newcommand{\fn}[1]{\texttt{\small #1}}

\begin{document}
\maketitle

\begin{abstract}
Se documenta el avance en la inversión cinemática del sismo intraslab de Calama
del 2020-06-03 ($M_w$ 6.8, $-23.247^\circ$, $-68.530^\circ$, profundidad
\SI{123.4}{\km}) con el paquete \fn{kdellipspy} (parche elíptico de
\emph{slip}, 7 parámetros, algoritmo Neighbourhood). El objetivo es resolver la
geometría de la aspereza, que las funciones de Green sintéticas 1D (AXITRA) no
logran constreñir pese a buenos ajustes de forma de onda. Se construyó un flujo
de función de Green empírica (EGF) y se mostró, mediante un control limpio, que
la ``aspereza compacta'' obtenida inicialmente era un artefacto de errores de
alineación temporal, no una característica robusta de los datos.
\end{abstract}

\section{Datos y preprocesamiento}
Registros de acelerómetros de las redes CSN/IPOC. Todos los canales (\fn{HN*} y
\fn{HL*}) son acelerómetros en este conjunto; se integran una vez a velocidad.
El observado se construye a partir de los \fn{.txt} crudos: \emph{detrend},
\emph{demean}, \emph{taper}, integración (\fn{cumtrapz}), filtro
\emph{bandpass} Butterworth de 4º orden \textbf{causal} (idéntico al aplicado a
las sintéticas), y remuestreo a una rejilla común. El misfit es el L2
normalizado global
\begin{equation}
  \chi^2 = \frac{\sum_{s}\sum_{c\in W}(o_{sc}-u_{sc})^2}
                {\sum_{s}\sum_{c\in W} o_{sc}^2},
\end{equation}
evaluado en \textbf{ventanas de \SI{20}{\s}}: P (radial $+$ vertical) desde
$t_P-\SI{1}{\s}$ y S (transversal) desde $t_S-\SI{1}{\s}$ (tiempos de
\texttt{taup}, iasp91). Las amplitudes relativas entre estaciones importan; un
factor de amplitud uniforme se cancela.

\section{Reajuste del caso AXITRA \emph{legacy}}
Se reprodujo el resultado de referencia en \fn{inversions/calama2020}:
$N_\text{pts}=128$, $\Delta t=\SI{1}{\s}$, tiempo de origen corregido a
\texttt{07:35:34}, ventana de misfit de \SI{20}{\s}, filtro causal, y
\num{3000} evaluaciones del NA. Banda \SIrange{0.02}{0.1}{\Hz}, 8 estaciones.

\begin{center}
\begin{tabular}{lc}
\toprule
Cantidad & Valor \\
\midrule
Misfit (mejor) & \num{0.204} \\
$M_w$ recuperado & \num{6.91} \\
\bottomrule
\end{tabular}
\end{center}

\noindent Buen ajuste de forma, pero el \emph{appraisal} confirma (5.ª vez)
que la geometría queda sin resolver: misfit bajo $\neq$ resolución.

\section{Espectros}
\paragraph{Amplitud (desplazamiento).} Para las 48 estaciones se calculó el
espectro de desplazamiento en el dominio de la frecuencia,
$|U(f)| = |A(f)|/(2\pi f)^2$, evitando la deriva de la doble integración
temporal; combinación vectorial de las 3 componentes. Salida:
\fn{output/espectros\_amplitud.pdf}.

\paragraph{Señal vs.\ ruido (metodología ISOLA).} Siguiendo Sokos \&
Zahradník: ventana de ruido (NTW) del mismo largo que la de señal (STW) y
terminando donde ésta comienza; $\mathrm{SNR}(f)$ como promedio por componente
de $S_c(f)/N_c(f)$; valor único $=$ promedio de la curva sobre la banda de
inversión. Curvas guardadas en \fn{output/snr\_curvas.csv}, figura en
\fn{output/espectros\_snr.pdf}. \textbf{Hallazgo:} el ruido microsísmico
oceánico tiene su máximo justo en \SIrange{0.1}{0.3}{\Hz} (la banda EGF
elegida), y el SNR de la EGF mejora $2$--$3\times$ al desplazarse a
\SIrange{0.15}{0.4}{\Hz}.

\begin{table}[h]
\centering
\caption{SNR promedio (ISOLA) de la EGF $M$4.9 por estación.}
\begin{tabular}{lcc}
\toprule
Estación & \SIrange{0.1}{0.3}{\Hz} & \SIrange{0.15}{0.4}{\Hz} \\
\midrule
PB06 & 7.1 & 24.5 \\
PB09 & 5.6 & 23.2 \\
AF01 & 4.5 & 17.9 \\
PB03 & 4.1 & 14.6 \\
PB19 & 4.7 & 12.0 \\
PB01 & 4.2 & 11.3 \\
PB05 & 3.8 & 10.0 \\
\bottomrule
\end{tabular}
\end{table}

\section{Función de Green empírica (EGF)}
\paragraph{Selección y validación.} Búsqueda FDSN (56 candidatos); el mejor es
el evento \textbf{2020-06-16 $M$4.9} ($\Delta_\text{3D}=\SI{17.7}{\km}$),
con CC$_S$(T) mediana $0.77$ en \SIrange{0.1}{0.3}{\Hz}. Lección: a baja
frecuencia pesa más el SNR (magnitud) que la co-localización fina.

\paragraph{Operador.} \emph{Shift-and-sum} con directividad,
\begin{equation}
  u_i(t) = \sum_k \frac{m_k}{M_0^\text{egf}}\,
           \mathrm{EGF}_i\!\left(t-\tau_{ki}\right), \qquad
  \tau_{ki} = t_\text{rup}(k) + \frac{r_{ki}-r_{0i}}{c},
\end{equation}
con retardo fraccional vía rampa de fase en FFT. Tres tests de cordura pasan
(límite puntual exacto, retardo, conservación de momento a baja frecuencia).
La geometría se toma de \fn{kdellipspy}
(\fn{build\_geometry\_with\_ellipse\_slip} $\to$ momentos y tiempos de ruptura).

\paragraph{Revalidación CC por banda.} La CC mainshock$\leftrightarrow$EGF se
sostiene en \SIrange{0.15}{0.4}{\Hz} (mediana $0.72$) y \textbf{recupera PB02}
(de $0.48$ a $0.73$): su problema en la banda baja era el microsismo, no el
sitio. Pasa de 7/8 a \textbf{8/8} estaciones válidas, ganando cobertura
azimutal. Los lags azimutales de directividad se preservan
(norte $+5/+6$ s, sur $-1$ s).

\section{Inversión EGF: el problema de la compacidad}
La primera batería (banda \SIrange{0.1}{0.3}{\Hz}, sin corrección de lags,
9 corridas) entregó por primera vez un paisaje discriminante (aleatorios
$\sim 1.12$ vs.\ mejor $0.87$) y una \emph{aspereza compacta} ($a_1\,a_2 \approx
3\times6$ km, área $40$--$80$ km$^2$). Sin embargo, el piso de misfit
$\sim0.87$ y los pisos de fase por ventana ($1-\mathrm{CC}^2$) sugerían que el
ajuste estaba limitado por desalineaciones estáticas de $2$--$4$ s
(localización del catálogo de la EGF $+$ trayecto 3D).

\paragraph{Lags estáticos.} Se añadió al misfit un \textbf{corrimiento estático
por estación y ventana} ($\pm\SI{4}{\s}$, P y S independientes), y se excluyeron
las ventanas P de menor SNR (PB02, PB09). Nueva batería de 6 corridas en
\SIrange{0.15}{0.4}{\Hz} con dos controles para aislar el efecto.

\begin{table}[h]
\centering
\caption{Mejor modelo por corrida (\num{3000} evaluaciones NA c/u).
Área $=\pi a_1 a_2$.}
\small
\begin{tabular}{lccccccccc}
\toprule
Corrida & banda & lags & misfit & $a_1$ & $a_2$ & dmax & $v_r$ & área & $M_w$ \\
        & (Hz)  &      &        & (km)  & (km)  & (m)  & (km/s)& (km$^2$) & \\
\midrule
principal           & 0.15--0.4 & $\pm4$ & 0.629 & 8.0 & 9.8  & 6.8 & 4.1 & 246 & 7.24 \\
semilla 22          & 0.15--0.4 & $\pm4$ & 0.556 & 8.5 & 12.3 & 7.4 & 1.8 & 330 & 7.35 \\
semilla 33          & 0.15--0.4 & $\pm4$ & 0.621 & 9.8 & 8.4  & 4.5 & 1.6 & 257 & 7.13 \\
sin PB02            & 0.15--0.4 & $\pm4$ & 0.601 & 6.6 & 12.0 & 7.7 & 1.9 & 249 & 7.28 \\
\textbf{nolag (ctrl)} & 0.15--0.4 & no    & 0.767 & 7.7 & \textbf{3.8} & 7.3 & 1.5 & \textbf{92} & 6.97 \\
\textbf{0.1--0.3+lag (ctrl)} & 0.1--0.3 & $\pm4$ & 0.543 & 5.2 & 11.6 & 7.7 & 1.6 & 189 & 7.21 \\
\bottomrule
\end{tabular}
\end{table}

\paragraph{Conclusión.} El control limpio (misma banda, estaciones y
exclusiones; sólo lags on/off) es inequívoco: \textbf{92 km$^2$ sin lags vs.\
246--330 km$^2$ con lags}. Los errores de alineación penalizaban las fuentes
extendidas y favorecían artificialmente la solución compacta. La compacidad de
la primera batería era, en gran medida, un \emph{artefacto de timing}.

\section{Estado actual y recomendación}
\begin{itemize}
  \item \textbf{Robusto:} fuente \emph{extendida} (área $\sim190$--330 km$^2$),
        dmax $\sim\SI{7}{\m}$; el paisaje discrimina ($\Delta\approx0.3$ sobre
        aleatorios).
  \item \textbf{No resuelto:} orientación ($\theta$), posición temporal ($t_p$),
        $v_r$ ($1.6$--$4.1$); $M_0$ \emph{sobreestimado} (factor 3--6,
        $M_w$ 7.0--7.35 vs.\ 6.8), acoplado a la incertidumbre de
        $M_0^\text{egf}$.
  \item \textbf{Riesgo metodológico:} lags libres de $\pm4$ s eliminan el
        \emph{moveout} azimutal, que es la señal que constriñe la fuente finita
        $\Rightarrow$ varios lags rielan en $+4$ s y la geometría se
        desordena.
\end{itemize}

\paragraph{Recomendación.} (i) Mantener AXITRA 1D
(\SIrange{0.02}{0.1}{\Hz} \emph{legacy}; \SIrange{0.06}{0.15}{\Hz} evento nuevo)
como resultado de producción para lo que sí constriñe ($M_0$, mecanismo,
profundidad, duración). (ii) Tratar la EGF como línea complementaria para la
geometría, pero \textbf{reemplazar los lags libres por una alineación estática
medida una sola vez} (lag de CC, \fn{cc\_bandas.csv}) horneada en los datos $+$
lag residual $\pm1$ s, preservando así la directividad. (iii) Calibrar
$M_0^\text{egf}$ por cociente espectral de \emph{plateaus} a baja frecuencia
antes de interpretar dmax$/M_0$ en términos absolutos.

\paragraph{Reproducibilidad.} Flujo en
\fn{inversions/2020-06-03/egf/codigos/}; inversión vía
\fn{invertir\_egf.py} con \fn{-{}-banda}, \fn{-{}-csv}, \fn{-{}-lagmax},
\fn{-{}-sinP}, \fn{-{}-dmax}. Informes: \fn{INFORME\_NOCHE.md} (1.ª batería),
\fn{INFORME\_b1540.md} (lags estáticos).

\end{document}
